\section{Repaso y complementos de la configuración Dense vs. MoE}

\subsection{La naturaleza de $d_{\text{model}}$}

$d_{\text{model}}$ (es decir, \texttt{hidden\_size}) es la dimensión central que atraviesa todo el modelo:
\begin{itemize}
  \item La dimensión del token embedding de entrada de cada capa es $d_{\text{model}}$
  \item La dimensión del token embedding de salida de cada capa sigue siendo $d_{\text{model}}$
  \item La consistencia de $d_{\text{model}}$ garantiza la viabilidad de la \textbf{conexión residual} (las dimensiones antes y después de la transformación son iguales, por lo que se pueden sumar directamente)
\end{itemize}

\subsection{Dimensión intermedia y estructura de la FFN}

\begin{table}[H]
\centering
\caption{Comparación de la dimensión intermedia de FFN/MoE}
\begin{tabular}{lcc}
\toprule
& \textbf{Dense (32B)} & \textbf{MoE (30B-A3B)} \\
\midrule
Dimensión intermedia $d_{\text{ff}}$ & 25600 ($5 \times d_{\text{model}}$) & 768 (por cada Expert) \\
Dirección del mapeo & $5120 \to 25600$ (up) & $2048 \to 768$ (down!) \\
Número de Experts & --- & 128 \\
Número de Experts activados & --- & 8 \\
\bottomrule
\end{tabular}
\end{table}

\begin{warningbox}{La dimensión intermedia del Expert de MoE no es de ``up''}
En el modelo Dense, la dimensión intermedia de la FFN es $25600 > d_{\text{model}} = 5120$, un proceso de «aumento de dimensión». Pero la dimensión intermedia de cada Expert en MoE es $768 < d_{\text{model}} = 2048$, lo que en realidad es un proceso de «reducción de dimensión». Aunque en el código se sigue usando el nombre $W_{\text{up}}$, en esencia cada Expert es más «delgado».
\end{warningbox}

\subsection{Configuración del Decoder Sparse Stack}

En la configuración de Qwen3-30B-A3B hay dos parámetros clave:
\begin{itemize}
  \item \texttt{decoder\_sparse\_step}: controla cada cuántas capas se usa una MLP Dense (cuando el residuo del módulo es 0, se usa MLP Dense)
  \item \texttt{mlp\_only\_layers}: especifica qué capas usan la MLP Dense estándar
\end{itemize}

En la configuración actual, \textbf{la parte de FFN de las 48 capas usa MoE}, sin que ninguna capa recurra a la MLP Dense.

\subsection{Resumen de la sección}

El modelo MoE se «reduce» en todas las dimensiones ($d_{\text{model}}$, número de cabezas, número de capas), y compensa la capacidad expresiva mediante 128 Experts independientes. La FFN de cada Expert es «angosta» (768 dimensiones), pero, al combinar varios Experts, el ancho equivalente puede llegar a 6144.

